\section{Invariants}\label{sec:invariants}

Theorem~\ref{thm:slice} states that each slice amplitude $a_u$ is a passive scalar carried by the divergence-free flow $\boldsymbol U_u=\uE+su\,\bhat$ of \eqref{eq:Uu}. The invariants of KRMHD derived in this section all follow from that statement. The hypotheses are those of the theorem: $\Lam\notin[0,1]$; $\Phi$ and $\Psi$ are smooth in a periodic box of volume $V$ and otherwise arbitrary, so that they may be forced, resistive or prescribed; and $G$ obeys the branch equation \eqref{eq:branch} with no collisions, dissipation or forcing. We write $\langle\cdot\rangle=V^{-1}\int\dd^3r\,(\cdot)$ for the volume average. The last subsection, on finite $\kperp\rho_i$, leaves KRMHD and states only what can be proved about the advection terms there.

\subsection{The quadratic family}\label{inv:sec-family}

For a real function $\varphi$, the operator $\varphi(L)$ is defined by the functional calculus of the self-adjoint operator $L$ (theorem~\ref{thm:spectrum}): the slice amplitude of $\varphi(L)G$ is $\varphi(u)\,a_u$. Since $G\mapsto a$ is unitary from $\mathcal W$ with the free-energy product \eqref{eq:fe} onto $L^2(\mathbb R,\dd u)$,
\begin{equation}\label{eq:Qphi}
\Qphi\equiv\int\dd^3r\,\fe{G}{\varphi(L)G}=\int\dd u\;\varphi(u)\int\dd^3r\,a_u^2=V\!\int\varphi(u)\,\mathcal E(u)\,\dd u,\qquad \mathcal E(u)\equiv\langle a_u^2\rangle .
\end{equation}
The function $\mathcal E(u)$ is the distribution of free energy over phase velocity. Its integral is $\langle\fe{G}{G}\rangle$, and $\tfrac12\mathcal E(u)\,\dd u$ is the free energy per unit volume, in units of $n_iT_i$, carried by the slices with phase velocities in $[u,u+\dd u]$. It is the diagonal $C(u,u)$ of the two-slice correlation studied in \S\,\ref{sec:decorrelation}.

\begin{proposition}\label{inv:energy}
Under the hypotheses of theorem~\ref{thm:slice}, with no collisions, dissipation or forcing acting on $G$, the slice energy $\mathcal E(u)$ is independent of time for every $u$ (for almost every $u$ if $G$ is not continuous in $u$). Hence $\Qphi$ is constant in time for every $\varphi$ for which the integral in \eqref{eq:Qphi} exists at $t=0$, either absolutely or as a principal value.
\end{proposition}

\begin{proof}
By \eqref{eq:slice}, $\partial_t a_u^2+\boldsymbol U_u\cdot\nabla a_u^2=0$. Because $\nabla\cdot\boldsymbol U_u=0$, this is $\partial_t a_u^2+\nabla\cdot(\boldsymbol U_u\,a_u^2)=0$, and integration over the periodic box gives $\dd\mathcal E(u)/\dd t=0$. The function $u\mapsto\mathcal E(u)$ is therefore the same at all times, and so is every integral of it.
\end{proof}

The family $\{\Qphi\}$ contains exactly the information in $\mathcal E(u)$. Its conservation states that, in the absence of collisions, free energy never moves between phase velocities; what changes in time is the correlation between different slices, which is the subject of \S\,\ref{sec:decorrelation}. A dissipation that acts identically on all velocities, such as $\kperp$ hyperdiffusion applied to every $g_m$, damps each slice separately \eqref{eq:sliceforcing}, so that $\mathcal E(u)$ decreases at every $u$ without exchange between slices. Collisions that depend on velocity couple the slices.

Whether a given $\Qphi$ is finite depends on $\varphi$ and on the state, and by proposition~\ref{inv:energy} it is enough to check this at $t=0$. For $\varphi(u)=u^n$ the integral converges absolutely if $\int\dd^3r\int u^{2k}G^2/F_M\,\dd u<\infty$ for some integer $k\ge n/2$. Indeed $|u|^n\le1+u^{2k}$, and $\int u^{2k}a_u^2\,\dd u=\fe{L^kG}{L^kG}$ is finite exactly when $G$ lies in the domain of $L^k$. That domain is the domain of $U^k$, because $L=UP$ \eqref{eq:LUP} and $P-I$ has rank one with range spanned by $F_M$, which has moments of all orders. For $\varphi(u)=1/u$ we define
\[
Q_{1/u}=\lim_{\delta\to0^+}V\!\int_{|u|>\delta}\frac{\mathcal E(u)}{u}\,\dd u .
\]
The limit exists when $\mathcal E$ is H\"older continuous near $u=0$. This is the case when $G(\boldsymbol r,\cdot)$ is H\"older continuous near $u=0$, uniformly in $\boldsymbol r$, because $a_u$ is a combination of $G$ and of its Hilbert transform with coefficients that are smooth in $u$ (\S\,\ref{sec:slices}), and the Hilbert transform preserves local H\"older continuity.

The conservation of the polynomial members can also be checked without the transform. Multiplying \eqref{eq:branch} by $u^k$ and integrating gives the moment hierarchy
\[
\big(\partial_t+\{\Phi,\cdot\}\big)M_k=-s\gradpar M_{k+1}+\frac{s\,c_{k+1}}{\Lam}\,\gradpar M_0,\qquad c_k=M_k[F_M],
\]
whose members $k=0$ and $k=1$ are equations (B23) and (B24) of \citet{Chandran2026}, hereafter CMM. Expanding the resolvent form \eqref{inv:resolvent} below at large $z$ gives $\fe{G}{L^nG}=\int u^nG^2/F_M\,\dd u-\Lam^{-1}\sum_{a+b+j=n}e_j\,M_aM_b$, where $1/\eps(z)=\sum_je_jz^{-j}$. With these relations $\dd Q_{u^n}/\dd t$ reduces to a combination of the antisymmetric integrals $\int\dd^3r\,M_a\gradpar M_b$ that vanishes identically; we checked this symbolically for $n\le8$. We also checked against GANDALF's right-hand side that the members of the truncated family of \S\,\ref{sec:truncated}, which include $W$ and $\Gamma$, are conserved to round-off.

Per-velocity quadratic invariants have precedents in simpler settings. \citet{PlunkParker2014} observed that, in the linear problem at a single wavenumber, any $u$-weighted integral of the squared transformed distribution is conserved; this is the linear form of $\Qphi$ for every $\varphi$. \citet{Plunk2015} proved a nonlinear per-velocity quadratic invariant for a quasi-two-dimensional drift-kinetic model without parallel streaming, in which the transform that diagonalises the linear dynamics does not depend on the wavenumber and therefore commutes with $\boldsymbol E\times\boldsymbol B$ advection. The transformed electrostatic drift-kinetic equation of \citet[][eq.~85]{Heninger2018}, for a three-dimensional slab with Boltzmann electrons, has every transformed component streaming along the uniform guide field and all components carried by one self-consistent $\boldsymbol E\times\boldsymbol B$ flow; it implies per-slice invariants as well, although they are not stated there. What these systems share with KRMHD is a Landau coupling that does not depend on the wavenumber, so that the transform commutes with the advection, and an advecting flow that is divergence free. Passivity is not required: in the drift-kinetic system the slices remain coupled through the self-consistent potential, and each still conserves its energy. \citet{Kunz2015} remarked that, with collisions ordered out of their pressure-anisotropic generalisation of KRMHD, the free-energy invariant they consider is one of an infinite number of invariants of the system; they did not specify the others.

\subsection{The invariants of CMM}\label{inv:sec-cmm}

CMM found three quadratic invariants for each KRMHD branch: the free energy $W^\pm$ (their B13), $\Gamma^\pm$ (B25) and the phase-space helicity $\Hph^\pm$ (B18). Written in the variable $u=v_\parallel/\vth$, with $G^\pm(v_\parallel)\,\dd v_\parallel=G^\pm(u)\,\dd u$, they are the members $\varphi=1$, $u$ and $1/u$ of \eqref{eq:Qphi}:
\begin{equation}\label{eq:WGH}
\begin{aligned}
W^\pm&=\frac{n_iT_i}{2}\,\big\langle\fe{G^\pm}{G^\pm}\big\rangle=\frac{n_iT_i}{2}\int\mathcal E^\pm(u)\,\dd u,\\
\Gamma^\pm&=\frac{\vth}{2}\,\big\langle\fe{G^\pm}{LG^\pm}\big\rangle=\frac{\vth}{2}\int u\,\mathcal E^\pm(u)\,\dd u,\\
\Hph^\pm&=\frac{n_iT_i}{2\vth}\,\Big\langle\PVint\frac{(G^\pm)^2}{uF_M}\,\dd u\Big\rangle=\frac{n_iT_i}{2\vth}\,\PVint\frac{\mathcal E^\pm(u)}{u}\,\dd u ,
\end{aligned}
\end{equation}
where $L$, the product $\fe{\cdot}{\cdot}$ and $\mathcal E^\pm$ are those of branch $\pm$, with $\Lam=\Lam^\pm$, and $\vth=s$ in units with $v_A=1$. The first expression in each line is CMM's definition rewritten in $u$; the second follows from \eqref{eq:Qphi}, and for $\Hph^\pm$ from \eqref{inv:pv} below. The normalisations are CMM's: $W^\pm$ carries the factor $n_iT_i$ and $\Gamma^\pm$ does not, so that $n_iT_i\Gamma^\pm$ has the dimensions of $\vth W^\pm$. CMM's free energy $\tilde W_{\rm compr}$ (B14) and helicity $\tilde H_{\textrm{ph-sp}}$ (B19) of the Laguerre sectors $l\ge2$ are the members $\varphi=1$ and $\varphi=1/u$ for $\Lam=\infty$, summed over $l$. In the Hermite representation $G=\sum_mg_mH_m(u)F_M/\sqrt{2^mm!}$ of \S\,\ref{sec:model},
\[
\fe{G}{G}=\sum_mp_m\,g_m^2=\sum_mg_m^2-\frac{g_0^2}{\Lam},\qquad
\tfrac12\fe{G}{LG}=\tfrac12\,g^{\mathsf T}\Pw A\,g=\frac{1}{\sqrt2}\Big(\sum_m\sqrt{m+1}\,g_mg_{m+1}-\frac{g_0g_1}{\Lam}\Big),
\]
which are CMM's (B31) and (B32), the latter divided by $\vth$. The helicity has no finite Hermite form of this kind; its truncation is the subject of \eqref{eq:Hoddeven} in \S\,\ref{sec:truncated}.

The kernel of $\Hph^\pm$ contains no $\Lam$. Formally this follows from $L=UP$. The operator $P$ is invertible for $\Lam\neq1$, with $P^{-1}=I+F_MM_0/(\Lam-1)$, so $PL^{-1}=PP^{-1}U^{-1}=U^{-1}$ and
\[
\fe{G}{L^{-1}G}=\ip{G}{PL^{-1}G}=\ip{G}{U^{-1}G}=\int\frac{G^2}{uF_M}\,\dd u .
\]
This is literal when $G/u$ lies in $\mathcal W$. For a general $G$, which need not vanish at $u=0$, the principal value follows from the resolvent form $R_G(z)=\fe{G}{(z-L)^{-1}G}$ with $\operatorname{Im}z\neq0$. By theorem~\ref{thm:spectrum}, $R_G$ is the Stieltjes transform of $a_u^2$. Writing $(z-L)X=G$ as $(z-u)X=G-uF_M\,M_0[X]/\Lam$ and taking $M_0$ gives $M_0[X]=C_G(z)/\eps(z)$, with $C_G(z)=\int G(v)/(z-v)\,\dd v$ and $\eps$ the dielectric function \eqref{eq:eps}. Substituting $X$ into $\fe{G}{X}$,
\begin{equation}\label{inv:resolvent}
R_G(z)=\int\frac{a_u^2}{z-u}\,\dd u=\int\frac{G^2}{(z-u)F_M}\,\dd u-\frac{z}{\Lam}\,\frac{C_G(z)^2}{\eps(z)} .
\end{equation}
Let $G$ be H\"older continuous near $u=0$, and let $z=\ii\eta$ with $\eta\to0^+$. By the Sokhotski--Plemelj formula the two integrals tend to $-\PVint a_u^2/u\,\dd u-\ii\pi a_0^2$ and to $-\PVint G^2/(uF_M)\,\dd u-\ii\pi G(0)^2/F_M(0)$, and the last term vanishes, because $C_G(\ii\eta)$ stays bounded and $\eps(0+\ii0)=1-1/\Lam\neq0$. Hence, at every point $\boldsymbol r$,
\begin{equation}\label{inv:pv}
\PVint\frac{a_u^2}{u}\,\dd u=\PVint\frac{G^2}{uF_M}\,\dd u,\qquad a_0^2=\frac{G(0)^2}{F_M(0)} .
\end{equation}
The helicity kernel, multiplication by $1/(uF_M)$, is therefore the same for both branches and for the Laguerre sectors, and $\Hph^\pm$ depends on $\Lam^\pm$ only through $G^\pm$. The kernels of $W^\pm$ and $\Gamma^\pm$ in the product $\ip{\cdot}{\cdot}$, which are $P$ and $PUP$, contain $\Lam$ through the rank-one part of $P$.

\subsection{The complex helicity and the slice at zero velocity}\label{inv:sec-complex}

CMM's gyrokinetic helicity (their 2.15) uses the weight $1/(v_\parallel-\ii0)$ rather than a principal value, and the imaginary part in their Plemelj split (2.17) is the contribution of $v_\parallel=0$. The analogue for one branch, defined like (B18) but with this weight, is the member $\varphi(u)=1/(u-\ii0)=\mathrm P(1/u)+\ii\pi\delta(u)$ of the family. By \eqref{eq:aunit} and \eqref{inv:pv},
\begin{equation}\label{eq:Hcomplex}
\Hph^\pm_{\rm c}\equiv\frac{n_iT_i}{2\vth}\Big\langle\int\frac{(G^\pm)^2}{(u-\ii0)F_M}\,\dd u\Big\rangle
=\Hph^\pm+\ii\,\frac{\pi n_iT_i}{2\vth}\,\big\langle\varrho(0)\,\Tr[G^\pm](0)^2\big\rangle
=\Hph^\pm+\ii\,\frac{\pi n_iT_i}{2\vth}\,\Big\langle\frac{G^\pm(\boldsymbol r,0,t)^2}{F_M(0)}\Big\rangle .
\end{equation}
Its imaginary part is proportional to $\mathcal E^\pm(0)$, the energy of the slice at $u=0$, which proposition~\ref{inv:energy} conserves. That slice is $G$ itself at $u=0$. Because $(LG)(0)=0$, the branch equation at $u=0$ reads
\[
\partial_tG(\boldsymbol r,0,t)+\{\Phi,G(\boldsymbol r,0,t)\}=0,
\]
and, since $1-1/\Lam>0$ for every $\Lam\notin[0,1]$, the slice amplitude is $a_0=G(\boldsymbol r,0,t)/F_M(0)^{1/2}$. The slice at zero phase velocity is advected by the $\boldsymbol E\times\boldsymbol B$ flow alone, with neither streaming nor the magnetic-flutter term. CMM showed that, in full gyrokinetics, $\psi_s=h_s-(q_s/T_{0s})\langle\chi\rangle_{\boldsymbol R_s}F_{0s}$ at $v_\parallel=0$ is advected by $\langle\chi\rangle_{\boldsymbol R_s}$ alone (their 2.22), which gives the conservation of ${\rm Im}\,\Hph$ at each $v_\perp$ (their 2.23). In KRMHD, where the streaming term of the kinetic equation (CMM's 3.25) vanishes at $v_\parallel=0$, their statement is the $u=0$ member of the slice family, and \eqref{eq:Hcomplex} is its form for one branch.

\subsection{Non-quadratic and cross-sector invariants}\label{inv:sec-casimir}

Each slice is a Lagrangian invariant of the volume-preserving flow $\boldsymbol U_u$: $a_u(\boldsymbol r,t)=a_u(\boldsymbol X_{u,t}^{-1}(\boldsymbol r),0)$, where $\boldsymbol X_{u,t}$ is the flow map of $\boldsymbol U_u$. Therefore, for every $u$ and every function $C$ for which the integral exists,
\begin{equation}\label{eq:casimir}
\frac{\dd}{\dd t}\int\dd^3r\,C(a_u)=0 .
\end{equation}
For smooth solutions this also follows from $C'(a)\{A,a\}=\{A,C(a)\}=\nabla\cdot[C(a)\,\zhat\times\nabla A]$ and $C'(a)\,\partial_za=\partial_zC(a)$, applied with $A=\Phi+su\Psi$. The quadratic member $C(a)=a^2$ gives $V\mathcal E(u)$. By \eqref{eq:sameflow}, the slice amplitudes $a^\pm_u$ of the two branches and $a^{(l)}_u=G_l(u)/F_M(u)^{1/2}$ of the Laguerre sectors $l\ge2$, for which $\Lam=\infty$, are carried by the same flow $\boldsymbol U_u$ and so share its flow map. Hence
\begin{equation}\label{eq:crosssector}
\frac{\dd}{\dd t}\int\dd^3r\,C\big(a^+_u,a^-_u,a^{(2)}_u,a^{(3)}_u,\dots\big)=0,\qquad\text{for example}\qquad\frac{\dd}{\dd t}\int\dd^3r\,a^+_u\,a^-_u=0 .
\end{equation}
Integrating the last relation over $u$ with any weight gives quadratic invariants that couple the two branches.

These statements hold for fields that are continuous in space. Under spatial Fourier--Galerkin truncation, as in pseudo-spectral codes, only the quadratic members survive, because the projected bracket conserves $\int a^2\,\dd^3r$ but in general not $\int C(a)\,\dd^3r$. A Hermite code that evolves both branches does not conserve the cross-branch forms exactly either, because the discrete phase velocities of the two branches, the eigenvalues of the ladder matrix at $\Lam^+$ and at $\Lam^-$, differ; the exception is the node $u=0$, which both branches have when $M$ is even (\S\,\ref{sec:truncated}).

In two dimensions ($\partial_z=0$), invariants labelled by velocity are known in several reduced models: Lagrangian invariants advected by velocity-dependent stream functions in the two-field reconnection models of \citet{Schep1994}, \citet{Cafaro1998} and \citet{Grasso2001} and in their drift-kinetic continuum \citep{Liseikina2004,Pegoraro2004}; per-velocity invariants of two-dimensional gyrokinetics \citep[][app.~F]{Schekochihin2009} and of the low-beta fluid-kinetic equations of \citet{ZoccoSchekochihin2011}; and the per-velocity quadratic invariants that survive Galerkin truncation, which \citet{ZhuHammett2010} used to construct absolute equilibria of two-dimensional gyrokinetics. In these models streaming along the field breaks the invariants in three dimensions. In three dimensions, \citet{Tassi2022} found that Hermite truncations of an electromagnetic drift-kinetic model with an odd number of moments contain one field that is advected by the $\boldsymbol E\times\boldsymbol B$ flow alone and has an infinite family of Casimirs; it corresponds to the slice at $u=0$ (\S\,\ref{sec:truncated}). The invariants \eqref{eq:casimir} at $u\neq0$ and the cross-sector invariants \eqref{eq:crosssector}, which hold in three dimensions with streaming along the perturbed field, are new.

\subsection{The local conservation law and \texorpdfstring{$\Gamma$}{Gamma} as a flux}\label{inv:sec-flux}

The free energy of each branch is conserved locally, with a flux that has a simple form.

\begin{proposition}\label{prop:localflux}
Let $G$ solve \eqref{eq:branch} with $\Phi$ and $\Psi$ smooth, and let $G$ and its first spatial derivatives lie in the domain of $U$ at every point. Then the free-energy density $w=\tfrac12\fe{G}{G}$, per unit $n_iT_i$, and the density $\gamma=\tfrac{s}{2}\fe{G}{LG}$ satisfy, at every point,
\begin{equation}\label{eq:localflux}
\partial_tw+\nabla\cdot\big(w\,\uE+\gamma\,\bhat\big)=0 .
\end{equation}
In slice variables $w=\tfrac12\int a_u^2\,\dd u$ and $\gamma=\tfrac{s}{2}\int u\,a_u^2\,\dd u$, so that $w\,\uE+\gamma\,\bhat=\tfrac12\int a_u^2\,\boldsymbol U_u\,\dd u$.
\end{proposition}

\begin{proof}
The product $\fe{\cdot}{\cdot}$ acts on $u$ at fixed $\boldsymbol r$ with a kernel that does not depend on $\boldsymbol r$, so every derivation $D$ in $\boldsymbol r$ satisfies $D\fe{f}{h}=\fe{Df}{h}+\fe{f}{Dh}$. Both $\{\Phi,\cdot\}$ and $\gradpar$ are such derivations, and both commute with $L$, which acts on $u$ only. Since $L$ is symmetric in $\fe{\cdot}{\cdot}$ at each point,
\[
\gradpar\fe{G}{LG}=\fe{\gradpar G}{LG}+\fe{G}{\gradpar LG}=\fe{\gradpar LG}{G}+\fe{G}{\gradpar LG}=2\,\fe{G}{\gradpar LG}.
\]
Taking $\fe{G}{\cdot}$ of \eqref{eq:branch} therefore gives $\partial_tw+\{\Phi,w\}+\gradpar\gamma=0$. Finally $\{\Phi,w\}=\uE\cdot\nabla w=\nabla\cdot(w\,\uE)$ and $\gradpar\gamma=\bhat\cdot\nabla\gamma=\nabla\cdot(\gamma\,\bhat)$, because $\nabla\cdot\uE=\nabla\cdot\bhat=0$. The slice forms follow from the unitarity of $G\mapsto a$ and from the fact that the slice amplitude of $LG$ is $u\,a_u$.
\end{proof}

The flux in \eqref{eq:localflux} is the sum over slices of the slice energy times the slice velocity. Averaged over the box, \eqref{eq:localflux} gives $\dd W^\pm/\dd t=0$, and the flux itself has a direct meaning. Since $\uE$ has no $z$ component and $\hat b_z=1$, the $z$ component of the volume-averaged flux is $\langle\gamma\rangle$, which by \eqref{eq:WGH} is $\Gamma^\pm$:
\begin{equation}\label{eq:Gammaflux}
\Gamma^\pm=\zhat\cdot\big\langle w\,\uE+\gamma\,\bhat\big\rangle,\qquad
\frac{n_iT_i\,\Gamma^\pm}{W^\pm}=\vth\,\frac{\int u\,\mathcal E^\pm(u)\,\dd u}{\int\mathcal E^\pm(u)\,\dd u}.
\end{equation}
Thus $n_iT_i\Gamma^\pm$ is the net flux of the free energy $W^\pm$ along $\boldsymbol B_0$: the $z$ component of the volume-averaged free-energy flux, or equivalently the flux per unit area through a plane $z={\rm const}$, averaged over $z$. The perpendicular components of $\langle\gamma\,\bhat\rangle=\langle\gamma\,\boldsymbol b_\perp\rangle$ need not vanish, and $\Gamma^\pm$ says nothing about them. The ratio $n_iT_i\Gamma^\pm/W^\pm$ is the energy-weighted mean of the parallel velocities $u\vth$ at which the slices stream, and $\Gamma^\pm=0$ when $\mathcal E^\pm(u)$ is even in $u$. CMM derived the conservation of $\Gamma^\pm$ from the moment equations (their B20--B26); its reading as a field-aligned flux of free energy is new.

\subsection{Finite \texorpdfstring{$\kperp\rho_i$}{k-perp rho-i}}\label{inv:sec-finite}

The results above use two properties of KRMHD: the kernel $P\varphi(L)$ of every member of the family is independent of the wavenumber, and the flow that advects $G$ does not depend on $v_\perp$. Both fail at finite $\kperp\rho_i$. In the kinetic equation of the isothermal-electron-fluid (ITEF) system (CMM's 3.16), the ion distribution is advected by the gyroaverage of the potential, which depends on $v_\perp$, and by the term $(\rho_i\vth/2)\{\hat v_\perp^2\hat J_1\delta B/B,\cdot\}$, where $\hat v_\perp=v_\perp/\vth$ and $\hat J_0$, $\hat J_1$, $\hat\Gamma_0$, $\hat\Gamma_1$ are CMM's gyroaveraging operators (their 3.11--3.13). The field relations acquire coefficients that depend on $\kperp$, so that the coupling that plays the role of $\Lam$ depends on $\boldsymbol k$, as in the gyrokinetic form of the transform of \citet[][eq.~79]{Heninger2018}; the invariant of \citet{Plunk2015} likewise rests on a transform that is independent of the wavenumber, which requires $k^2\rho^2\ll1$. What can be proved concerns the advection terms alone.

Let $\theta(\boldsymbol r,\boldsymbol v,t)=\sum_{\boldsymbol k}\theta_{\boldsymbol k}(\boldsymbol v)\,e^{\ii\boldsymbol k\cdot\boldsymbol r}$ be real and advected by a prescribed real potential $\chi(\boldsymbol r,t)$, so that $\partial_t\theta_{\boldsymbol k}=\sum_{\boldsymbol p+\boldsymbol q=\boldsymbol k}(\boldsymbol p\times\boldsymbol q)_z\,\chi_{\boldsymbol p}M_{\boldsymbol p}\theta_{\boldsymbol q}$. Here $M_{\boldsymbol p}=I$ if $\chi$ does not depend on velocity, which is the equation $\partial_t\theta+\{\chi,\theta\}=0$, and $M_{\boldsymbol p}$ is multiplication by $J_0(p_\perp v_\perp/\Omega_i)$ if each velocity is advected by the gyroaveraged potential. Consider the quadratic forms $Q=\sum_{\boldsymbol k}\theta_{\boldsymbol k}^\dagger K_{\boldsymbol k}\theta_{\boldsymbol k}$, with $K_{\boldsymbol k}$ Hermitian operators on functions of velocity and $K_{-\boldsymbol k}=\bar K_{\boldsymbol k}$. Writing $\dd Q/\dd t=2\,{\rm Re}\sum_{\boldsymbol k}\theta_{\boldsymbol k}^\dagger K_{\boldsymbol k}\partial_t\theta_{\boldsymbol k}$ and relabelling half of the terms by $(\boldsymbol p,\boldsymbol q)\to(-\boldsymbol p,\boldsymbol q+\boldsymbol p)$ gives the identity
\begin{equation}\label{inv:triad}
\frac{\dd Q}{\dd t}=\sum_{\boldsymbol p,\boldsymbol q}(\boldsymbol p\times\boldsymbol q)_z\,{\rm Re}\Big[\chi_{\boldsymbol p}\,\theta_{\boldsymbol q+\boldsymbol p}^\dagger\big(K_{\boldsymbol q+\boldsymbol p}M_{\boldsymbol p}-M_{\boldsymbol p}K_{\boldsymbol q}\big)\theta_{\boldsymbol q}\Big].
\end{equation}
For $\theta$ supported on $\pm\boldsymbol q$, $\pm(\boldsymbol q+\boldsymbol p)$ and $\chi$ on $\pm\boldsymbol p$, with $\boldsymbol p$ and $\boldsymbol q$ in general position, the sum has four equal terms and $\dd Q/\dd t=4(\boldsymbol p\times\boldsymbol q)_z\,{\rm Re}[\chi_{\boldsymbol p}\theta_{\boldsymbol q+\boldsymbol p}^\dagger(K_{\boldsymbol q+\boldsymbol p}M_{\boldsymbol p}-M_{\boldsymbol p}K_{\boldsymbol q})\theta_{\boldsymbol q}]$.

\begin{lemma}[Advection independent of velocity]\label{inv:lemA}
Let $M_{\boldsymbol p}=I$. Then $Q$ is conserved for every $\chi$ and every $\theta$ if and only if $K_{\boldsymbol k}$ is the same for all $\boldsymbol k$ with $\boldsymbol k_\perp\neq0$. The $K_{\boldsymbol k}$ with $\boldsymbol k_\perp=0$ are unconstrained.
\end{lemma}

\begin{proof}
The coefficient $(\boldsymbol p\times\boldsymbol q)_z$ vanishes when $\boldsymbol q_\perp=0$ and when $(\boldsymbol q+\boldsymbol p)_\perp=0$, so every term of \eqref{inv:triad} that survives involves two wavevectors with non-zero perpendicular parts; if $K_{\boldsymbol k}$ is the same for all of these, every term vanishes. Conversely, if $K_{\boldsymbol k_1}\neq K_{\boldsymbol k_2}$ with $\boldsymbol k_{1\perp}$ and $\boldsymbol k_{2\perp}$ non-zero and not parallel, the triad with $\boldsymbol q=\boldsymbol k_1$ and $\boldsymbol p=\boldsymbol k_2-\boldsymbol k_1$ has $(\boldsymbol p\times\boldsymbol q)_z\neq0$, and $\theta_{\boldsymbol q}$, $\theta_{\boldsymbol q+\boldsymbol p}$ and $\chi_{\boldsymbol p}$ can be chosen to make $\dd Q/\dd t$ non-zero. Wavevectors with parallel perpendicular parts are compared through a third wavevector.
\end{proof}

\begin{lemma}[Gyroaveraged advection]\label{inv:lemAp}
Let $M_{\boldsymbol p}$ be multiplication by $J_0(p_\perp v_\perp/\Omega_i)$, and let $K_{\boldsymbol k}=K$ for all $\boldsymbol k$. Then $Q$ is conserved for every $\chi$ and every $\theta$ if and only if $K$ commutes with every $M_{\boldsymbol p}$. This holds when $K$ acts separately at each $v_\perp$, whatever it does in $v_\parallel$, and fails when $K$ couples different values of $v_\perp$.
\end{lemma}

\begin{proof}
With $K_{\boldsymbol k}=K$ the operator in \eqref{inv:triad} is the commutator $[K,M_{\boldsymbol p}]$. If it vanishes for every $\boldsymbol p$, so does $\dd Q/\dd t$; if it is non-zero for some $\boldsymbol p$, which then has $\boldsymbol p_\perp\neq0$, a triad with $(\boldsymbol p\times\boldsymbol q)_z\neq0$ and suitable $\theta$ makes $\dd Q/\dd t$ non-zero.
\end{proof}

\begin{lemma}[Helicity kernel]\label{inv:lemB}
Let $L_{\boldsymbol k}=UP_{\boldsymbol k}$ with $P_{\boldsymbol k}=I-F_MM_0/\Lam_{\boldsymbol k}$ and $\Lam_{\boldsymbol k}\notin[0,1]$ depending on $\boldsymbol k$. Then $P_{\boldsymbol k}L_{\boldsymbol k}^{-1}=U^{-1}$ on the range of $U$, for every $\boldsymbol k$. The form $\fe{G}{L_{\boldsymbol k}^{-1}G}_{\boldsymbol k}=\int G^2/(uF_M)\,\dd u$, taken as a principal value as in \eqref{inv:pv}, therefore does not depend on $\boldsymbol k$, and its kernel is a multiplication operator, whereas the kernels $P_{\boldsymbol k}$ and $P_{\boldsymbol k}UP_{\boldsymbol k}$ of $Q_1$ and $Q_u$ depend on $\boldsymbol k$ through $\Lam_{\boldsymbol k}$.
\end{lemma}

\begin{proof}
$P_{\boldsymbol k}$ is invertible for $\Lam_{\boldsymbol k}\neq1$, with $P_{\boldsymbol k}^{-1}=I+F_MM_0/(\Lam_{\boldsymbol k}-1)$, and maps the domain of $U$ onto itself. Hence $L_{\boldsymbol k}^{-1}=P_{\boldsymbol k}^{-1}U^{-1}$ on the range of $U$, and $P_{\boldsymbol k}L_{\boldsymbol k}^{-1}=U^{-1}$.
\end{proof}

Lemmas~\ref{inv:lemA}--\ref{inv:lemB} concern advection by a prescribed potential. They identify two obstructions that the advection terms place on a quadratic invariant at finite $\kperp\rho_i$: a kernel that depends on $\boldsymbol k$, as $P_{\boldsymbol k}$ and $P_{\boldsymbol k}UP_{\boldsymbol k}$ would through a coupling that depends on $\boldsymbol k$, and a kernel that couples different values of $v_\perp$. The branch forms $W^\pm$, $\Gamma^\pm$ and $\Hph^\pm$ have the second property, because they act on $G^\pm$, which are weighted $v_\perp$-integrals of $g$ (CMM's B2, B3 and B5). The helicity escapes the first obstruction by lemma~\ref{inv:lemB}. CMM's gyrokinetic helicity, whose kernel $1/(v_\parallel F_{0s})$ acts on $\psi_s$ and is a multiplication operator in velocity, also escapes the second, and CMM's proof shows that the remaining terms cancel. The lemmas do not give a necessary condition for invariance under the full equations, because terms that come from the field equations can make a form invariant even when its kernel depends on $\boldsymbol k$ and couples different velocities. The gyrokinetic free energy (CMM's 2.13; \citealt{Schekochihin2009}) is such a form when it is written in terms of the gyrocentre distribution. We checked this mechanism in a two-dimensional electrostatic gyrokinetic model with Boltzmann electrons: the field energy, whose kernel depends on $\boldsymbol k$ and couples velocities, is conserved to round-off by the self-consistent nonlinearity, and it is not conserved when the same advecting potential is prescribed.

What happens to $\Gamma^\pm$ at finite $\kperp\rho_i$ is open. CMM report $W$ and $\Hph$ as the quadratic invariants of general gyrokinetics; $\Gamma^\pm$ appears in their work only for KRMHD, and they do not state that no further invariants exist. Under the KRMHD ordering, in which the normalised potential $2\Phi/(\rho_i\vth)$ exceeds $\delta n/n$ and $\delta B/B$ by a factor $(\kperp\rho_i)^{-1}$ \citep[][CMM \S\,3.2]{Schekochihin2009}, the ITEF equations contain five terms of relative order $\kperp\rho_i$ that are absent from KRMHD: the advection of the ion distribution by $(\rho_i\vth/2)\,\hat v_\perp^2\hat J_1\delta B/B$ in the ion kinetic equation (CMM's 3.16); the term $(Z\vth/2\tau)\gradpar(\delta n/n)$ in the equation for the parallel vector potential (3.6); the bracket $(\rho_i\vth/2)\{(Z/\tau)\delta n/n,\delta B/B\}$ in the electron continuity equation (3.7); and the terms in which $1-\hat\Gamma_0$ and $1-\hat\Gamma_1$ act on the potential, in quasineutrality (3.10) and in perpendicular pressure balance (3.14). The corrections that $\hat J_0$ and $\hat J_1$ make to the gyroaverages are of order $\kperp^2\rho_i^2$ (CMM's 3.11). Whether some correction of $\Gamma^\pm$ at order $\kperp\rho_i$ is conserved in the presence of these terms has not been computed. In the two-dimensional electrostatic model mentioned above, where quasineutrality makes the coupling $\Lam_{\boldsymbol k}$ depend on $\kperp$ and $L_{\boldsymbol k}=UP_{\boldsymbol k}$, the analogue of $\Gamma$, $\sum_{\boldsymbol k}\fe{g}{L_{\boldsymbol k}g}_{\boldsymbol k}$, is conserved at $\kperp\rho_i=0$ and is broken by the nonlinearity at a rate proportional to $(\kperp\rho_i)^2$ (log--log slope 1.99 over a factor of eight in $\rho_i$). That model lacks the electromagnetic terms of order $\kperp\rho_i$ listed above and does not settle the question for gyrokinetics. Related work at finite Larmor radius is that of \citet{CostelloPlunk2025}, who revisit the quasi-two-dimensional invariant of \citet{Plunk2015} and extend it to finite Larmor radius.
